\section{Formal Causal Normalization}
\label{app:causality}
\subsection{Hypotheses}
Fix an address graph $\mathcal F$ with unit-edge shortest-path distance and a family of positive-depth admitted histories. Every primitive event has address increment zero or one, preserves the declared event structure, and is recoverable in the complete transducer description. The depth of a history is the irreducible number of such primitive steps in the chosen grammar.

Under these assumptions the triangle inequality proves $d_{\mathrm{field}}(s,t)\leq|\gamma|_{\mathrm{event}}$. If at least one depth-one address-changing event is admitted, the supremum in \eqref{eq:cth} is exactly one. The first clause is the bound; the existence clause supplies saturation. No statement about all long geodesics follows from the one-step witness.

\subsection{An illustrative saturating family}
As an example independent of the Thalean carrier, take an unbounded line of identical reversible relays with a fixed admitted payload transfer to the next address. A payload can pass $n$ adjacent relationships in $n$ steps. Its endpoint distance is $n$, so the causal slope is one at every length. If the transfer is already injective, a fixed or purely operational receipt can suffice.

This example demonstrates consistency of the causal normalization with an unbounded field. It is not a derivation of the Thalean field wiring, a photon model, or a physical lattice. A native field construction must prove its own payload and residue compatibility.

\subsection{Why no clock calibration is hidden}
One can introduce an external clock for a simulator and assign a duration to an update. That produces a rate in simulation units, not a native physical speed. A physical identification would require the same apparatus to establish an event-to-clock relation, a field-distance-to-length relation, and invariance of the resulting observable speed under the intended frame transformations.

The locked definition remains a dimensionless causal datum. It should be revised only if a native theorem, contradiction, or failed audit forces a change, not because a different denominator gives a more familiar-looking formula.
